\documentclass[12pt,a4paper]{article}
\usepackage[utf-8]{inputenc}
\usepackage[vietnamese]{babel}
\usepackage{geometry}
\usepackage{tabularx}
\usepackage{booktabs}
\usepackage{array}
\usepackage{xcolor}
\usepackage[hidelinks]{hyperref}

\geometry{margin=1in}
\setlength{\parindent}{0pt}
\setlength{\parskip}{0.5cm}

\title{Bảng Kiểm Thử Hệ Thống Quản Lý Nhà Trọ}
\author{Room Minh House Management System}
\date{\today}

\begin{document}

\maketitle

\tableofcontents
\newpage

\section*{Hướng Dẫn Sử Dụng Bảng Kiểm Thử}

Tài liệu này chứa các trường hợp kiểm thử cho hệ thống quản lý nhà trọ, được tổ chức theo từng module chức năng:

\begin{itemize}
    \item Quản lý Phòng
    \item Quản lý Người Thuê
    \item Quản lý Hợp Đồng
    \item Quản lý Hóa Đơn
    \item Quản lý Công Tơ Điện Nước
\end{itemize}

Mỗi bảng kiểm thử bao gồm các cột:
\begin{itemize}
    \item \textbf{STT}: Số thứ tự của trường hợp kiểm thử
    \item \textbf{Chức năng}: Tên chức năng hoặc tính năng cần kiểm thử
    \item \textbf{Dữ liệu nhập}: Dữ liệu đầu vào hoặc điều kiện tiên quyết
    \item \textbf{Kết quả mong đợi}: Kết quả dự kiến khi thực hiện chức năng
    \item \textbf{Kết quả}: Kết quả thực tế (để lại trống để điền khi kiểm thử)
\end{itemize}

\newpage

% Placeholder for tables - will be added in subsequent tasks

\end{document}
